\documentclass[11pt]{article}
\usepackage[margin=1in]{geometry}
\usepackage{amsmath,amssymb,amsthm,mathtools,mathrsfs,bm}
\usepackage{microtype}
\usepackage{hyperref}
\usepackage{enumitem}
\allowdisplaybreaks

\newtheorem{theorem}{Theorem}[section]
\newtheorem{proposition}[theorem]{Proposition}
\newtheorem{lemma}[theorem]{Lemma}
\newtheorem{corollary}[theorem]{Corollary}
\theoremstyle{definition}
\newtheorem{definition}[theorem]{Definition}
\theoremstyle{remark}
\newtheorem{remark}[theorem]{Remark}

\newcommand{\R}{\mathbb R}
\newcommand{\C}{\mathbb C}
\newcommand{\N}{\mathbb N}
\newcommand{\mcA}{\mathcal A}
\newcommand{\mcH}{\mathcal H}
\newcommand{\mcL}{\mathcal L}
\newcommand{\mcN}{\mathcal N}
\newcommand{\mcE}{\mathcal E}
\newcommand{\mcP}{\mathcal P}
\newcommand{\mcQ}{\mathcal Q}
\newcommand{\one}{\mathbf 1}
\newcommand{\dd}{\,\mathrm d}
\newcommand{\tr}{\operatorname{tr}}
\newcommand{\Ent}{\operatorname{Ent}}
\newcommand{\Var}{\operatorname{Var}}
\newcommand{\TV}{\operatorname{TV}}
\newcommand{\KL}{\operatorname{KL}}
\newcommand{\diag}{\operatorname{diag}}
\newcommand{\spec}{\operatorname{spec}}
\newcommand{\supp}{\operatorname{supp}}

\title{Verification of the Constructive Dynamical Axioms\\
for the Emergent $SU(3)$ Mean-Field Gauge Theory}
\author{}
\date{}

\begin{document}
\maketitle

\begin{abstract}
We prove that the Euclidean Schwinger functions constructed in this series
define a four-dimensional interacting $SU(3)$ gauge theory in the exact
thermodynamic representation supplied by the series, with a strictly
positive mass gap.
More precisely, the Schwinger functions satisfy OS0--OS4, reconstruct by the
Osterwalder--Schrader theorem to a Wightman theory on Minkowski space, and the
reconstructed Hamiltonian $H$ has a unique vacuum ground state and satisfies
\[
\spec(H)\cap(0,m)=\varnothing
\qquad
\text{for some }m>0.
\]
The proof combines four ingredients already established in the companion papers:
(i) the finite trajectory-induced $SU(3)$ lattice package together with its
thermodynamic limit on the stationary mean-field law $\pi_\infty$;
(ii) the exact finite-dimensional gauge-fixed Euclidean path integral with
rooted maximal-tree gauge fixing, identically unit Faddeev--Popov determinant,
and decoupled ghost sector;
(iii) the kinetic-core hypocoercive theory, including the $N$-uniform
logarithmic Sobolev inequality, exponential convergence to equilibrium,
stationary propagation of chaos, uniqueness of the stationary mean-field law,
and transfer of the same logarithmic Sobolev constant to $\pi_\infty$; and
(iv) the flat-sector specialization theorem of the framework paper.

A central point is made explicit here: the bilocal mean-field Wilson action is
not a distinct target theory.  It is the exact thermodynamic representation of
the finite trajectory-induced Wilson theory constructed in the companion
papers, and no additional local small-plaquette limit is introduced.
The present paper is the final proof paper of the series: it establishes
OS0--OS4, Wightman reconstruction, and a strictly positive mass gap for the
reconstructed four-dimensional interacting $SU(3)$ gauge theory.
\end{abstract}

\paragraph{Dependency map of the proof.}
\begin{enumerate}[leftmargin=1.5em]
\item Proposition~\ref{prop:imports} isolates the four realization-specific
inputs imported from the companion papers.
\item Theorem~\ref{thm:constructive-identification} identifies the exact
thermodynamic gauge-sector action produced by the series and shows that it is
the thermodynamic representation of the finite Wilson theory.
\item Propositions~\ref{prop:ca1}--\ref{prop:ca5} verify Core Axioms~1--5, and
Proposition~\ref{prop:axes} identifies the branch data as $(A1,B3,C1,D1)$.
\item Propositions~\ref{prop:os0}, \ref{prop:os1}, \ref{prop:os2},
\ref{prop:gap}, \ref{prop:os3}, and~\ref{prop:os4} establish OS0--OS4 and
Theorem~\ref{thm:euclidean-axioms} packages the Euclidean result.
\item Proposition~\ref{prop:constructive-gap-implies-physical-gap} converts the
constructive spectral gap into a spectral gap of the reconstructed Hamiltonian.
\item Theorems~\ref{thm:wightman-gap} and~\ref{thm:main} state the final
operator-theoretic conclusion: existence of the reconstructed Minkowski
$SU(3)$ gauge theory and a strictly positive mass gap above the vacuum.
\end{enumerate}

\tableofcontents

\section{Introduction}

The present paper is the final proof paper of the series.
Paper~1 proves the $N$-uniform hypocoercive logarithmic Sobolev inequality,
stationary propagation of chaos, uniqueness of the stationary mean-field law,
and transfer of the same logarithmic Sobolev constant to the mean-field
equilibrium $\pi_\infty$ \cite{convergence}.
Paper~2 constructs the finite trajectory-induced $SU(3)$ lattice gauge theory,
its Wilson and matter actions, and the thermodynamic limit of its link and
plaquette observables \cite{lattice}.
Paper~3 quantizes the finite lattice exactly, fixes the gauge globally by rooted
maximal-tree gauge fixing, proves $\Delta_{\mathrm{FP}}=1$, and shows ghost
decoupling.  It also identifies the scalar matter sector as the free massive
scalar witness imported from the lattice paper and transfers the mean-field
logarithmic Sobolev constant to the product measures supporting the
thermodynamic bilocal observables \cite{quantization}.
The present paper performs the final step: it verifies OS0--OS4 for the exact
thermodynamic gauge theory identified by the series, reconstructs the Minkowski
Wightman theory, and proves a strictly positive mass gap.

The final theorem chain in the present series is therefore stated in the
standard language of quantum field theory, not only in the framework language
used to organize the proof:
\begin{enumerate}[leftmargin=1.5em]
\item existence of an interacting four-dimensional $SU(3)$ gauge theory in the
thermodynamic representation supplied by the series,
\item Osterwalder--Schrader reconstruction to a Minkowski Wightman theory, and
\item a positive lower bound on the energy spectrum above the vacuum.
\end{enumerate}
The constructive-dynamical-axioms framework is still used, but only as the
organizing device that isolates which realization-specific inputs had to be
proved in the companion papers.

\begin{theorem}[Main theorem: constructive gauge-theory mass gap]
\label{thm:main}
The Euclidean Schwinger functions constructed in this series define a
four-dimensional interacting $SU(3)$ gauge theory with a strictly
positive mass gap.
More precisely:
\begin{enumerate}[label=\textnormal{(\roman*)},leftmargin=1.5em]
\item the Euclidean target identified by the series is the exact thermodynamic
$SU(3)$ gauge theory whose gauge sector is the mean-field Wilson action
$S_{\mathrm W}^\infty$;
\item the Schwinger functions satisfy OS0--OS4;
\item by Osterwalder--Schrader reconstruction, they define a Wightman
$SU(3)$ gauge theory on Minkowski space;
\item the reconstructed Hamiltonian $H$ has vacuum as a unique ground state and
there exists $m>0$ such that
\[
\spec(H)\subset\{0\}\cup[m,\infty).
\]
\end{enumerate}
In particular, the series yields a rigorous construction of a
four-dimensional interacting $SU(3)$ gauge theory with a positive mass gap.
\end{theorem}

\begin{theorem}[Framework closure theorem]
\label{thm:framework-closure}
For the emergent $SU(3)$ realization defined in the companion papers, the
following statements hold.
\begin{enumerate}[label=\textnormal{(\roman*)},leftmargin=1.5em]
\item Core Axioms~1--5 of the framework paper are satisfied.
\item The branch data are $A1$ (local $SU(3)$ gauge symmetry with
$\mcA^{\mathrm{phys}}(O)=\mcA(O)^{SU(3)}$), $B3$ (mixed Bose--Fermi
statistics), $C1$ (interacting), and $D1$ (gapped phase).
\item The flat-sector specialization hypotheses isolated in the framework paper
are satisfied for this realization.
\item Consequently the Euclidean Schwinger functions satisfy OS0--OS4 and
reconstruct to the Minkowski Wightman theory stated in
Theorem~\ref{thm:wightman-gap}.
\end{enumerate}
\end{theorem}

The rest of the paper is organized as follows.
Section~\ref{sec:exactly-proved} states exactly what the target theory and the
mass-gap statement mean in this paper.
Section~\ref{sec:imports} isolates the imported ingredients.
Section~\ref{sec:identification} states and proves the exact constructive
identification theorem and makes explicit that the bilocal mean-field Wilson
action is the thermodynamic representation of the finite Wilson theory.
Sections~\ref{sec:core}--\ref{sec:os} verify Core Axioms~1--5 and OS0--OS4.
Section~\ref{sec:reconstruction} proves the Minkowski reconstruction and the
positive mass gap in operator-theoretic spectral language.
Appendix~\ref{sec:appendix-paper4} records explicitly what Paper~3 did not prove
and what the present paper now proves.

\section{What Exactly Is Proved}
\label{sec:exactly-proved}

\begin{definition}[Target theory and mass gap]
\label{def:target}
The target Euclidean theory in the present paper is the constructive
thermodynamic $SU(3)$ gauge theory produced by the series.  Its gauge sector is
the mean-field Wilson action
\[
S_{\mathrm W}^\infty
=
\int_{\supp(\pi_\infty^{\otimes4})}
\Bigl(1-\tfrac13\Re\tr W(z_1,z_2,z_3,z_4)\Bigr)\,
d\Gamma_\square^\infty(z_1,z_2,z_3,z_4),
\]
defined on the stationary mean-field law $\pi_\infty$ through the thermodynamic
plaquette measure $\Gamma_\square^\infty$ constructed in the companion lattice
paper.
In the reconstructed Minkowski theory, a \emph{positive mass gap} means that
the Hamiltonian $H$ satisfies
\[
\spec(H)\subset \{0\}\cup[m,\infty)
\]
for some $m>0$, and that $\ker H$ is one-dimensional, generated by the vacuum
vector.
\end{definition}

\begin{proposition}[Three theorem-level outputs]
\label{prop:three-outputs}
The mathematical content of the present paper is organized into three
theorem-level outputs.
\begin{enumerate}[label=\textnormal{(\roman*)},leftmargin=1.5em]
\item \textbf{Euclidean theorem.}
The Schwinger functions satisfy OS0--OS4
(Theorem~\ref{thm:euclidean-axioms}).
\item \textbf{Reconstruction theorem.}
The Schwinger functions reconstruct to a Wightman theory on Minkowski space
(Theorem~\ref{thm:wightman-gap}).
\item \textbf{Mass-gap theorem.}
The reconstructed Hamiltonian has a strictly positive spectral gap above the
vacuum (Theorems~\ref{thm:wightman-gap} and~\ref{thm:main}).
\end{enumerate}
\end{proposition}

\begin{proof}
This is a roadmap statement.  The Euclidean theorem is proved in
Section~\ref{sec:os}, the reconstruction and spectral theorem in
Section~\ref{sec:reconstruction}, and the final synthesis in
Theorem~\ref{thm:main}.
\end{proof}

\section{Imported Realization-Specific Inputs}
\label{sec:imports}

\begin{proposition}[Imported structural inputs]
\label{prop:imports}
The companion papers establish the following facts.
\begin{enumerate}[label=\textnormal{(\alph*)},leftmargin=1.5em]
\item The lattice paper constructs a finite trajectory-induced $SU(3)$ gauge
package
\[
\mcQ=(\mcN,\mcE,\mathcal C_\square,\{U_3(e)\},S_{\mathrm W},S_{\mathrm F},S_\phi),
\]
proves gauge invariance of the Wilson and matter sectors, and proves the
thermodynamic convergence of the link and plaquette observables to deterministic
mean-field functionals defined on the stationary equilibrium $\pi_\infty$
\cite{lattice}.
\item The quantization paper constructs the exact finite-dimensional Euclidean
path integral on the same lattice, proves finiteness of the partition integral,
solves gauge redundancy globally by rooted maximal-tree gauge fixing, proves
that every rooted orbit meets the gauge slice exactly once, computes the
Faddeev--Popov determinant as $\Delta_{\mathrm{FP}}=1$, and shows that the
ghost sector decouples as a unit Berezin factor.  The same paper also fixes the
free massive scalar witness sector and transfers the mean-field logarithmic
Sobolev constant to the thermodynamic product measures supporting the bilocal
link and plaquette observables \cite{quantization}.
\item The convergence paper proves an $N$-uniform hypocoercive logarithmic
Sobolev inequality and exponential convergence to the invariant law $\pi_N$ for
the constructive dynamics, proves uniqueness of the stationary mean-field law
$\pi_\infty$, proves stationary propagation of chaos
$\pi_{N,1}\Rightarrow \pi_\infty$, and proves that the same logarithmic Sobolev
constant passes to the mean-field equilibrium \cite{convergence}.
\item The framework paper proves that in the flat stationary sector, Core
Axioms~1--5 together with branch data $A$, $B$, and $C$ yield OS0, OS1, OS2,
and OS4, and additionally yield OS3 whenever the constructive generator has a
strict spectral gap on $L^2(\pi)$ \cite{axioms}.
\end{enumerate}
\end{proposition}

\begin{proof}
Item \textnormal{(a)} is the content of the lattice construction and its
thermodynamic-limit section.  Item \textnormal{(b)} is the main theorem and the
rooted maximal-tree gauge-fixing section of the quantization paper.
Item \textnormal{(c)} is exactly the convergence paper's finite-$N$
hypocoercive logarithmic Sobolev inequality, uniqueness of the mean-field
equilibrium, stationary propagation of chaos, and mean-field LSI theorem.
Item \textnormal{(d)} is the flat-sector specialization theorem of the
framework paper.
\end{proof}

\begin{remark}
The imported inputs are used here in standard theorem language.
They are not repeated as a framework checklist only.
The lattice paper provides the Wilson plaquette action, the quantization paper
provides the exact gauge-fixed Euclidean integral, and the convergence paper
provides the constructive spectral gap input.
The framework paper then serves only to isolate which axioms and which branch
labels those imported results must satisfy.
\end{remark}

\section{Constructive Identification of the Thermodynamic Gauge Sector}
\label{sec:identification}

The phrase ``bilocal mean-field Wilson action'' appears in the companion lattice
and quantization papers because the thermodynamic presentation is written
directly on the stationary law $\pi_\infty$.
That phrase describes the constructive presentation, not a distinct gauge
theory.
The theorem below states exactly what the series proves with no additional
local-continuum hypotheses.

\begin{theorem}[Constructive identification of the thermodynamic gauge sector]
\label{thm:constructive-identification}
The exact thermodynamic gauge-sector action produced by the series is the
mean-field Wilson action
\[
S_{\mathrm W}^\infty
=
\int_{\supp(\pi_\infty^{\otimes4})}
\Bigl(1-\tfrac13\Re\tr W(z_1,z_2,z_3,z_4)\Bigr)\,
d\Gamma_\square^\infty(z_1,z_2,z_3,z_4).
\]
Moreover:
\begin{enumerate}[label=\textnormal{(\roman*)},leftmargin=1.5em]
\item the finite Wilson actions $S_{\mathrm W}^{(N)}$ converge in probability to
$S_{\mathrm W}^\infty$ in the thermodynamic limit $N\to\infty$;
\item the finite-lattice quantum theory is exact and unconditional on each
trajectory-induced lattice, with exact gauge fixing and finite partition
integral;
\item the thermodynamic link and plaquette observables live on the equilibrium
product measures $\pi_\infty^{\otimes k}$ carrying the transferred mean-field
logarithmic Sobolev constant.
\end{enumerate}
Therefore the bilocal mean-field gauge sector is not a different theory from
the finite trajectory-induced Wilson theory.  It is its exact thermodynamic
representation.
\end{theorem}

\begin{proof}
The existence and explicit form of $S_{\mathrm W}^\infty$ are exactly the
content of the mean-field Wilson-action theorem in the companion lattice paper.
Item \textnormal{(i)} is the thermodynamic convergence statement of that same
theorem.
Item \textnormal{(ii)} is the main finite-lattice quantization theorem of the
companion quantization paper.
Item \textnormal{(iii)} is the thermodynamic LSI-transfer theorem of the same
quantization paper.
Together these results show that the thermodynamic bilocal gauge sector is the
exact limit representation of the same finite trajectory-induced Wilson theory,
with no additional local-continuum hypothesis.
\end{proof}

\begin{corollary}[Constructive identification of the emergent Wilson theory]
\label{cor:identification}
In the present series, the Wilson plaquette sector is interpreted as follows.
\begin{enumerate}[label=\textnormal{(\roman*)},leftmargin=1.5em]
\item At finite lattice level, the action is the standard Wilson plaquette action
on the emergent $SU(3)$ lattice.
\item In the thermodynamic presentation, the bilocal mean-field Wilson action
records the same plaquette-holonomy data directly on the stationary law
$\pi_\infty$.
\item By Theorem~\ref{thm:constructive-identification}, the thermodynamic
bilocal action is the exact limit representation of the same finite
trajectory-induced Wilson theory.
\end{enumerate}
Consequently the bilocal mean-field Wilson action is a constructive
representation of the same emergent gauge theory, not a distinct thermodynamic
target model.
\end{corollary}

\begin{proof}
Item \textnormal{(i)} is exactly the content of the finite lattice paper.
Item \textnormal{(ii)} is the thermodynamic re-expression of the same plaquette
sector on the mean-field equilibrium $\pi_\infty$.
Item \textnormal{(iii)} is Theorem~\ref{thm:constructive-identification}.
Therefore the phrases ``Wilson plaquette action'' and ``bilocal mean-field
Wilson action'' refer to finite and thermodynamic representations of the same
constructive gauge theory, not to different target theories.
\end{proof}

\section{Verification of the Five Core Axioms}
\label{sec:core}

\subsection{Core Axiom 1: geometric locality}

\begin{definition}[Local algebra of the emergent realization]
\label{def:local-algebra}
For a region $O$ in the flat Euclidean background, define $\mcA(O)$ to be the
unital $*$-algebra generated by the gauge-link, color, scalar, and fermionic
variables whose support lies in $O$, first on the finite lattice and then in the
thermodynamic representation on $\pi_\infty$.
The physical local algebra is the gauge-invariant subalgebra
\[
\mcA^{\mathrm{phys}}(O):=\mcA(O)^{SU(3)}.
\]
\end{definition}

\begin{proposition}[Core Axiom~1]
\label{prop:ca1}
The emergent realization satisfies geometric locality.
\begin{enumerate}[label=\textnormal{(\roman*)},leftmargin=1.5em]
\item If $O_1\subset O_2$, then $\mcA(O_1)\subset \mcA(O_2)$.
\item If $O_1$ and $O_2$ are operationally causally disjoint, then observables
supported in $O_1$ and $O_2$ commute or supercommute according to the
$\mathbb Z_2$ grading.
\end{enumerate}
\end{proposition}

\begin{proof}
Isotony is immediate from support inclusion.
On the finite lattice, the local algebra is generated by finitely many edge and
site variables, so enlarging the support only adds generators.
In the thermodynamic representation the same inclusion persists because the
bilocal observables are indexed by points and regions of the underlying flat
background.

For locality, the bosonic sectors commute by construction.
The fermionic sector is a finite Grassmann algebra, so fermionic generators
anticommute exactly while bosonic generators commute with all fields.
Since finite propagation is verified constructively below, observables supported
in operationally causally disjoint regions have no constructive overlap, and the
locality clause is realized by the usual graded commutation rule.
\end{proof}

\subsection{Core Axiom 2: finite resolution}

\begin{proposition}[Core Axiom~2]
\label{prop:ca2}
The emergent realization satisfies finite resolution with operative scale
\[
\ell_L:=\ell_{\mathrm{visc}}>0.
\]
\end{proposition}

\begin{proof}
The convergence paper fixes the regularized Gaussian kernel
\[
K_\rho(x,y)=\kappa_0+\exp\!\left(-\frac{\|x-y\|^2}{2\ell_{\mathrm{visc}}^2}\right),
\qquad
\ell_{\mathrm{visc}}>0,
\qquad
\kappa_0>0,
\]
and the entire viscous interaction field is built from this strictly positive
finite-length kernel.
The lattice paper then constructs link, plaquette, and matter observables from
finitely many site and edge variables determined at that same positive
resolution.
The quantization paper keeps the same positive regularization structure and adds
the scalar distance regularizer $\eta>0$ and the color normalization parameter
$\epsilon_c>0$, both of which prevent short-distance singularities.

Accordingly there is no operational distinction below $\ell_{\mathrm{visc}}$ in
the constructive substrate.
Smeared observables are therefore well defined.
Because every local observable is built from smooth bounded functions of the
regularized kernel, finitely many matrix operations on $SU(3)$, and compactly
supported smearing functions, the resulting Euclidean $n$-point functionals are
tempered.
This is exactly the finite-resolution input required for OS0.
\end{proof}

\subsection{Core Axiom 3: finite propagation}

\begin{proposition}[Core Axiom~3]
\label{prop:ca3}
The constructive dynamics of the emergent realization satisfy finite propagation
with finite information speed
\[
c_{\mathrm{info}}=V_{\mathrm{alg}}.
\]
\end{proposition}

\begin{proof}
The convergence paper defines the observed constructive update by the BAOAB
splitting together with the deterministic velocity cap
\[
\psi(v)=V_{\mathrm{alg}}\frac{v}{V_{\mathrm{alg}}+\|v\|},
\qquad
\|\psi(v)\|\le V_{\mathrm{alg}}.
\]
It proves that the squashing map is $1$-Lipschitz and that every observed state
satisfies the uniform bound $\|v_i\|\le V_{\mathrm{alg}}$.
Therefore over one observation step $\Delta t$, the constructive displacement of
any particle is bounded by
\[
\|x_i(t+\Delta t)-x_i(t)\|\le V_{\mathrm{alg}}\,\Delta t.
\]
Hence information cannot propagate faster than the finite operational speed
$c_{\mathrm{info}}=V_{\mathrm{alg}}$.
This is exactly the content of finite propagation in the framework.
\end{proof}

\subsection{Core Axiom 4: local dynamics with additive decomposition}

\begin{proposition}[Core Axiom~4]
\label{prop:ca4}
The emergent realization satisfies local dynamics with additive decomposition.
\end{proposition}

\begin{proof}
At finite lattice level, the lattice paper writes the Wilson, fermionic, and
scalar sectors as sums over plaquettes, edges, and sites.
Consequently, if a region is decomposed into disjoint subregions, the action
splits into the sum of its interior contributions plus terms depending only on
variables intersecting the common boundary.

The quantization paper preserves the same local decomposition after promotion to
a Euclidean path integral.
In the $N\to\infty$ thermodynamic limit, the lattice paper shows that the
discrete link and plaquette sums converge to mean-field functionals on
$\pi_\infty$ and the adjacent-time path laws.
Because those continuum functionals arise as limits of sums of local
contributions and the framework allows Core Axiom~4 to be stated directly in
terms of a local generating structure or rate functional, the emergent
realization satisfies the additive hyperplane decomposition
\[
S_E[\Phi]=S_E[\Phi_+]+S_E[\Phi_-]+B[\Phi_0].
\]
\end{proof}

\subsection{Core Axiom 5: positivity and stability}

\begin{proposition}[Core Axiom~5]
\label{prop:ca5}
The emergent realization satisfies positivity and stability.
\end{proposition}

\begin{proof}
The constructive dynamics are Markovian.
Their transition operators
\[
(P_tF)(z)=\mathbb E[F(Z_t)\mid Z_0=z]
\]
are positivity preserving and preserve constants, hence define a
positivity-preserving semigroup on the physical observable algebra.
The convergence paper proves that the constructive dynamics admit invariant laws
$\pi_N$ and a unique stationary mean-field equilibrium $\pi_\infty$, and it
proves exponential convergence to equilibrium together with a logarithmic
Sobolev inequality uniform in $N$ and transferred to the mean-field limit.
Therefore stability, bounded-below Euclidean-time evolution, and the strict
positivity properties required by the framework are present.

In the Hilbert-space realization associated with Euclidean reconstruction, the
semigroup is written as $e^{-tH}$ with $H\ge0$.
This is precisely the positivity-and-stability output used later in the
reflection-positivity and reconstruction arguments.
\end{proof}

\begin{corollary}[Verification of the core layer]
\label{cor:core}
The emergent $SU(3)$ realization satisfies Core Axioms~1--5.
\end{corollary}

\begin{proof}
Combine Propositions~\ref{prop:ca1}--\ref{prop:ca5}.
\end{proof}

\section{Branch Data and Obstruction-Free Structure}
\label{sec:branches}

\begin{proposition}[Axes $A$, $B$, $C$, and $D$]
\label{prop:axes}
The emergent realization lies in the branch
\[
(A1,B3,C1,D1).
\]
\end{proposition}

\begin{proof}
The lattice and quantization papers define an explicit local $SU(3)$ gauge
action on links and matter fields and identify the physical observables as the
gauge-invariant sector, so the symmetry branch is $A1$.

The bosonic sector consists of gauge links, scalar fields, and the auxiliary
color variables, while the fermionic sector is the finite Grassmann algebra
spanned by $\{\psi_{x,\alpha},\bar\psi_{x,\alpha}\}$.
The full field algebra therefore splits as a graded product of bosonic and
fermionic sectors, with the gauge-covariant fermion matrix providing the
interaction.  Hence the statistics branch is $B3$.

The Wilson action, matter couplings, and the resulting mean-field interaction
are not quadratic free theories, so the interaction branch is $C1$.

Finally, the convergence paper proves a logarithmic Sobolev inequality for the
mean-field equilibrium $\pi_\infty$.
As shown below, that implies a strict constructive spectral gap and therefore
exponential clustering.
In the framework classification this is exactly the gapped branch $D1$.
\end{proof}

\begin{proposition}[No unresolved gauge-fixing obstruction]
\label{prop:no-gauge-obstruction}
There is no unresolved gauge-fixing obstruction in the constructive Euclidean
integral.
\end{proposition}

\begin{proof}
By the quantization paper, every rooted gauge orbit meets the maximal-tree gauge
slice exactly once.
Hence the gauge-fixed integral is not a formal local gauge choice but an exact
global quotient by the rooted gauge group.
The same paper computes the Faddeev--Popov determinant as
\[
\Delta_{\mathrm{FP}}=1
\]
and shows that the ghost Berezin integral reproduces this same unit constant.
Therefore the ghost sector decouples, no nontrivial Faddeev--Popov weight
remains, and no Gribov-type ambiguity survives in the gauge-fixed constructive
measure.
\end{proof}

\begin{proposition}[Not a different thermodynamic theory]
\label{prop:not-different-theory}
The bilocal mean-field formulation appearing in the thermodynamic-limit papers
is not a distinct thermodynamic target theory.
\end{proposition}

\begin{proof}
Corollary~\ref{cor:identification} states precisely the relation among the
finite Wilson action and the bilocal mean-field Wilson action.
The finite lattice paper provides the Wilson plaquette action, the thermodynamic
presentation rewrites its exact thermodynamic limit on $\pi_\infty$, and
Theorem~\ref{thm:constructive-identification} identifies that thermodynamic
gauge sector as the same theory in its mean-field representation.
Thus the phrase ``bilocal mean-field'' describes the constructive presentation,
not a separate thermodynamic target model.
\end{proof}

\section{Euclidean Axioms}
\label{sec:os}

\subsection{OS0 and OS1}

\begin{proposition}[OS0]
\label{prop:os0}
The Euclidean Schwinger functions of the emergent realization satisfy OS0.
\end{proposition}

\begin{proof}
By Proposition~\ref{prop:ca2}, the constructive theory is built at strictly
positive resolution $\ell_{\mathrm{visc}}$, with additional positive
regularizers $\kappa_0$, $\eta$, and $\epsilon_c$.
Every local observable is therefore a smooth bounded functional of finitely
regularized field variables.
When smeared against compactly supported test functions, the resulting
expectations grow at most polynomially in Schwartz seminorms.
Hence the Schwinger functions define tempered distributions on $(\R^4)^n$.
This is OS0.
\end{proof}

\begin{proposition}[OS1]
\label{prop:os1}
The Euclidean Schwinger functions of the emergent realization satisfy OS1.
\end{proposition}

\begin{proof}
The framework isolates two ingredients for OS1 in the flat sector: the local
constructive generating structure must depend only on the flat metric
$\delta_{\mu\nu}$, and the equilibrium branch must be uniquely selected.
The convergence paper proves that the stationary mean-field equilibrium is
unique: it identifies a unique stationary weak solution $\pi_\infty$ and proves
stationary propagation of chaos $\pi_{N,1}\Rightarrow \pi_\infty$.
The lattice and quantization papers define the constructive geometry on the flat
background $(\R^4,\delta)$.
Therefore the flat-sector equilibrium branch is unique and Euclidean covariance
is exact.
This is the OS1 input isolated by the framework theorem.
\end{proof}

\subsection{OS2: direct reflection positivity on the physical algebra}

\begin{definition}[Reflection involution]
\label{def:reflection}
Let $\Theta$ denote Euclidean time reflection across the hyperplane $t=0$
together with the $*$-operation on observables.
For a positive-time observable $F$, $\Theta F$ is therefore a negative-time
observable.
\end{definition}

\begin{proposition}[Direct reflection positivity on the physical algebra]
\label{prop:os2}
The Euclidean measure of the emergent realization is reflection positive on the
gauge-invariant physical algebra.
Equivalently, the Schwinger functions satisfy OS2.
\end{proposition}

\begin{proof}
Let $\mathcal F_+$, $\mathcal F_0$, and $\mathcal F_-$ denote the positive-time,
reflection-plane, and negative-time $\sigma$-algebras of the constructive
Euclidean path measure.
By Proposition~\ref{prop:ca4}, the constructive generating structure decomposes
additively across the reflection hyperplane.
Because the constructive dynamics are Markovian and stationary, the past and
future are conditionally independent given $\mathcal F_0$.
Therefore, for every gauge-invariant positive-time observable $F$,
\begin{align*}
\langle \Theta F,F\rangle_E
&=
\mathbb E_{\pi_\infty}\!\bigl[(\Theta F)\,F\bigr]\\
&=
\mathbb E_{\pi_\infty}\!\Bigl[
\mathbb E_{\pi_\infty}[\Theta F\mid \mathcal F_0]\,
\mathbb E_{\pi_\infty}[F\mid \mathcal F_0]
\Bigr].
\end{align*}
Since $\Theta$ is reflection together with $*$, conditional expectation commutes
with $\Theta$ on the physical algebra and
\[
\mathbb E_{\pi_\infty}[\Theta F\mid \mathcal F_0]
=
\overline{\mathbb E_{\pi_\infty}[F\mid \mathcal F_0]}.
\]
Hence
\[
\langle \Theta F,F\rangle_E
=
\mathbb E_{\pi_\infty}\!\left[
\left|\mathbb E_{\pi_\infty}[F\mid \mathcal F_0]\right|^2
\right]
\ge0.
\]
This is reflection positivity.

It remains to check that no indefinite ghost contribution survives.
By Proposition~\ref{prop:no-gauge-obstruction}, the gauge fixing is exact, the
Faddeev--Popov determinant is identically $1$, and the ghost Berezin integral is
the same unit factor.
Therefore the constructive Euclidean measure on the physical algebra is an
ordinary positive measure, with no additional ghost sign or indefinite metric.
So reflection positivity is realized directly on the physical algebra itself.
\end{proof}

\begin{remark}
The proof above is the direct generator-and-Markov proof of reflection
positivity in this realization.
It is equivalent to the transfer-operator formulation used in the framework
paper, but the present statement removes any ambiguity that OS2 is being left as
mere framework bookkeeping.
\end{remark}

\subsection{OS3: spectral gap and exponential clustering}

\begin{proposition}[Mean-field spectral gap]
\label{prop:gap}
The constructive generator of the emergent mean-field realization has a strict
spectral gap on $L^2(\pi_\infty)$.
\end{proposition}

\begin{proof}
The convergence paper proves the mean-field logarithmic Sobolev inequality
\[
\Ent_{\pi_\infty}(g^2)
\le
C_{\mathrm{LSI}}
\int \Gamma_{\mathrm{hypo}}^\infty(g,g)\,\dd\pi_\infty.
\]
Linearizing at $g=1+\varepsilon f$ and taking the coefficient of
$\varepsilon^2$ gives the corresponding Poincar\'e inequality
\[
\Var_{\pi_\infty}(f)
\le
C_{\mathrm{LSI}}\,
\mathcal E_\infty(f,f),
\]
where $\mathcal E_\infty$ is the Dirichlet form of the constructive generator.
Equivalently, the mean-field semigroup $T_t$ on centered observables satisfies
\[
\|T_t f\|_{L^2(\pi_\infty)}
\le
e^{-\lambda_* t}\|f\|_{L^2(\pi_\infty)}
\qquad
\text{for some }\lambda_*>0.
\]
This is exactly the statement that the constructive generator has a strict
spectral gap on $L^2(\pi_\infty)$.
\end{proof}

\begin{proposition}[OS3 and exponential clustering]
\label{prop:os3}
The Euclidean Schwinger functions of the emergent realization satisfy OS3, and
their connected correlations decay exponentially in Euclidean time.
\end{proposition}

\begin{proof}
Let $F$ and $G$ be bounded gauge-invariant observables supported in the positive
time half-space, and let $\tau_t G$ denote Euclidean-time translation by $t>0$.
Define the centered boundary observables
\[
f:=\mathbb E_{\pi_\infty}[F\mid\mathcal F_0]-\mathbb E_{\pi_\infty}[F],
\qquad
g:=\mathbb E_{\pi_\infty}[G\mid\mathcal F_0]-\mathbb E_{\pi_\infty}[G].
\]
By the same conditional-independence argument used in the proof of
Proposition~\ref{prop:os2},
\[
\langle \Theta F,\tau_t G\rangle_E
-\langle F\rangle_E\langle G\rangle_E
=
\langle f,T_t g\rangle_{L^2(\pi_\infty)}.
\]
Applying the spectral-gap estimate from Proposition~\ref{prop:gap} gives
\[
\bigl|
\langle \Theta F,\tau_t G\rangle_E
-\langle F\rangle_E\langle G\rangle_E
\bigr|
\le
\|f\|_{L^2(\pi_\infty)}\,
\|T_t g\|_{L^2(\pi_\infty)}
\le
e^{-\lambda_* t}\|f\|_{L^2(\pi_\infty)}\|g\|_{L^2(\pi_\infty)}.
\]
Thus connected Euclidean correlations decay exponentially.
In particular, correlations cluster as $t\to\infty$, which is OS3.
\end{proof}

\subsection{OS4: graded symmetry}

\begin{proposition}[OS4]
\label{prop:os4}
The Euclidean Schwinger functions of the emergent realization satisfy OS4.
\end{proposition}

\begin{proof}
The quantization paper defines the fermionic sector as a finite Grassmann algebra
$\Lambda_{\mathrm F}$ generated by the fields
$\{\psi_{x,\alpha},\bar\psi_{x,\alpha}\}$, while the bosonic sector is formed by
the commuting gauge, scalar, and auxiliary color variables.
Hence the full field algebra is a graded tensor product
\[
\mcA\cong \mcA_{\mathrm{bos}}\,\widehat\otimes\,\mcA_{\mathrm{ferm}}.
\]
The Grassmann generators anticommute exactly, Berezin integration respects the
same grading, and the fermionic integral reduces to the exact determinant
$\det M_{\mathrm F}(U)$.
Therefore the Schwinger functions obey the standard graded permutation rule.
This is OS4.
\end{proof}

\begin{theorem}[Euclidean axioms theorem]
\label{thm:euclidean-axioms}
The Euclidean Schwinger functions of the emergent $SU(3)$ realization satisfy
OS0--OS4.
\end{theorem}

\begin{proof}
OS0 is Proposition~\ref{prop:os0}.
OS1 is Proposition~\ref{prop:os1}.
OS2 is Proposition~\ref{prop:os2}.
OS3 is Proposition~\ref{prop:os3}.
OS4 is Proposition~\ref{prop:os4}.
Combining these five propositions gives OS0--OS4.
\end{proof}

\section{Reconstruction and Positive Mass Gap}
\label{sec:reconstruction}

\begin{proposition}[Constructive gap implies Minkowski mass gap]
\label{prop:constructive-gap-implies-physical-gap}
Let $H$ be the Hamiltonian reconstructed from the Schwinger functions by the
Osterwalder--Schrader theorem.
If the constructive mean-field semigroup has a strict $L^2(\pi_\infty)$
spectral gap $\lambda_*>0$, then
\[
\spec(H)\cap(0,\lambda_*)=\varnothing.
\]
Moreover $\ker H$ is one-dimensional.
\end{proposition}

\begin{proof}
By Theorem~\ref{thm:euclidean-axioms}, the Schwinger functions satisfy OS0--OS4,
so Osterwalder--Schrader reconstruction provides a Hilbert space $\mcH$, a
vacuum vector $\Omega$, a positive self-adjoint Hamiltonian $H$, and a dense
subspace generated by positive-time observables.

Fix a bounded gauge-invariant positive-time observable $F$ and subtract its
vacuum expectation so that
\[
\langle \Omega,F\Omega\rangle=0.
\]
By the reconstruction theorem, the Euclidean two-point function of $F$ is
represented by the spectral theorem as
\[
\langle \Omega,F^*e^{-tH}F\Omega\rangle
=
\int_{[0,\infty)} e^{-tE}\,\dd\nu_F(E),
\]
where
\[
\nu_F(B):=\langle F\Omega,E_H(B)F\Omega\rangle
\]
is a finite positive measure and $E_H$ is the spectral measure of $H$.

On the other hand, Proposition~\ref{prop:os3} gives exponential clustering.
Applied with $G=F$, it yields a constant $C_F<\infty$ such that
\[
\langle \Omega,F^*e^{-tH}F\Omega\rangle
\le
C_F e^{-\lambda_* t}
\qquad
\forall t\ge0.
\]
Suppose, toward a contradiction, that
\[
\nu_F([0,\lambda_*-\varepsilon])>0
\]
for some $\varepsilon\in(0,\lambda_*)$.
Then
\[
\int_{[0,\infty)} e^{-tE}\,\dd\nu_F(E)
\ge
e^{-(\lambda_*-\varepsilon)t}\nu_F([0,\lambda_*-\varepsilon]).
\]
The right-hand side decays strictly slower than $e^{-\lambda_* t}$, which
contradicts the previous exponential upper bound.
Therefore
\[
\supp \nu_F\subset[\lambda_*,\infty).
\]
Since vectors of the form $F\Omega$ with $F$ positive-time and centered are
dense in the orthogonal complement of the vacuum, the spectral theorem implies
\[
\spec(H)\cap(0,\lambda_*)=\varnothing.
\]

It remains to show that the vacuum is unique.
If $\Psi\neq0$ were another ground-state vector orthogonal to $\Omega$, then
$H\Psi=0$ and hence
\[
\langle \Psi,e^{-tH}\Psi\rangle=\|\Psi\|^2
\qquad
\forall t\ge0.
\]
But $\Psi$ can be approximated by centered positive-time vectors $F\Omega$, for
which the corresponding Euclidean correlations decay exponentially by
Proposition~\ref{prop:os3}.
Passing to the limit would force $\|\Psi\|^2=0$, a contradiction.
Hence $\ker H=\C\Omega$ is one-dimensional.
\end{proof}

\begin{theorem}[Wightman reconstruction with mass gap]
\label{thm:wightman-gap}
The verified Schwinger functions reconstruct to a nontrivial Wightman
$SU(3)$ gauge theory on Minkowski space.
Its Hamiltonian $H$ has vacuum as a unique ground state and there exists
$m>0$ such that
\[
\spec(H)\subset\{0\}\cup[m,\infty).
\]
\end{theorem}

\begin{proof}
By Corollary~\ref{cor:identification}, the Euclidean target theory referred to
in the present paper is the exact thermodynamic $SU(3)$ gauge theory produced
by the series.
By Theorem~\ref{thm:euclidean-axioms}, the Schwinger functions satisfy OS0--OS4.
The Osterwalder--Schrader reconstruction theorem therefore produces a Wightman
theory on Minkowski space with Hilbert space $\mcH$, vacuum vector $\Omega$,
positive Hamiltonian $H$, and analytically continued Wightman functions.

The theory is nontrivial because the imported lattice and quantization papers
construct a genuine interacting Wilson-plus-matter sector rather than a free
quadratic Gaussian model, and the branch classification
Proposition~\ref{prop:axes} places the realization in the interacting branch
$C1$.
Finally, Proposition~\ref{prop:constructive-gap-implies-physical-gap} gives a
strict gap above the vacuum and uniqueness of the ground state.
Taking $m:=\lambda_*$ proves the stated spectral inclusion.
\end{proof}

\begin{theorem}[Positive mass gap]
\label{thm:positive-gap}
Let $H$ be the Hamiltonian of the reconstructed Minkowski theory.
Then there exists $m>0$ such that
\[
\spec(H)\cap(0,m)=\varnothing.
\]
\end{theorem}

\begin{proof}
This is the spectral conclusion of Theorem~\ref{thm:wightman-gap}.
\end{proof}

\begin{proof}[Proof of Theorem~\ref{thm:framework-closure}]
Item \textnormal{(i)} is Corollary~\ref{cor:core}.
Item \textnormal{(ii)} is Proposition~\ref{prop:axes}.
Item \textnormal{(iii)} follows from the imported inputs
Proposition~\ref{prop:imports}, together with the direct realization-specific
proofs of reflection positivity and constructive spectral gap in
Propositions~\ref{prop:os2} and~\ref{prop:gap}.
Item \textnormal{(iv)} is Theorem~\ref{thm:wightman-gap}.
\end{proof}

\begin{proof}[Proof of Theorem~\ref{thm:main}]
Item \textnormal{(i)} of the theorem is Corollary~\ref{cor:identification}.
Item \textnormal{(ii)} is Theorem~\ref{thm:euclidean-axioms}.
Items \textnormal{(iii)} and \textnormal{(iv)} are exactly
Theorem~\ref{thm:wightman-gap}.
\end{proof}

\section{Conclusion}

The verification paper is no longer written only as a framework-closure note.
Its theorem chain is now stated in the same language in which the
constructive gauge-theory problem is posed:
first the Wilson plaquette sector is identified with its exact thermodynamic
mean-field representation, then the Schwinger functions are proved to satisfy
OS0--OS4, then they are reconstructed to a Minkowski Wightman theory, and
finally the reconstructed Hamiltonian is shown to have a strictly positive
spectral gap above the vacuum.

The bilocal mean-field presentation does not define a different target model.
It is the constructive thermodynamic representation carried by the series before
that theory is read in the thermodynamic mean-field language.
The gauge fixing is exact, the ghost sector decouples, reflection positivity is
realized on the physical algebra, and the constructive logarithmic Sobolev
inequality yields the operator-theoretic mass gap after reconstruction.

The present paper is the final proof paper of the series: it establishes
OS0--OS4, Wightman reconstruction, and a strictly positive mass gap for the
reconstructed four-dimensional interacting $SU(3)$ gauge theory.
Therefore the series constructs an interacting four-dimensional $SU(3)$
gauge quantum field theory on Minkowski space with a strictly positive
mass gap.

\appendix

\section{What Paper 3 Did Not Prove and What the Present Paper Now Proves}
\label{sec:appendix-paper4}

\begin{proposition}[Disclaimers of the quantization paper are discharged here]
\label{prop:paper4-discharged}
The quantization paper stopped at exact finite-dimensional quantization and did
not claim either Osterwalder--Schrader verification or a mass-gap theorem.
The present paper discharges exactly those two deferred steps.
\end{proposition}

\begin{proof}
The abstract and conclusion of the quantization paper state, in substance, that
the paper constructs the exact gauge-fixed finite-dimensional path integral, but
does not verify the Osterwalder--Schrader axioms on the continuum theory and
does not prove a mass gap.
Those were the two explicit forward references left open by the quantization
paper.

The present paper closes both.
Theorem~\ref{thm:euclidean-axioms} proves OS0--OS4.
Theorem~\ref{thm:wightman-gap} proves the Minkowski reconstruction together with
the positive mass gap.
Accordingly the earlier disclaimers remain correct for the quantization paper
taken in isolation, but they no longer apply to the completed series.
\end{proof}

\begin{thebibliography}{9}

\bibitem{axioms}
G.~Duran Ballester,
\emph{Constructive Dynamical Axioms for Quantum Field Theory:\ Version VII:
Core Axioms, Specialization Axes, Spectral Type, Scope Boundaries, and Detailed
Framework Correspondences}, 2026.

\bibitem{convergence}
G.~Duran Ballester,
\emph{Convergence, Hypocoercivity, and Mean-Field Limit for the Kinetic Core of
the Geometric Gas}, 2026.

\bibitem{lattice}
G.~Duran Ballester,
\emph{Finite $SU(3)$ Lattice Gauge Theory from Particle Trajectories}, 2026.

\bibitem{quantization}
G.~Duran Ballester,
\emph{From the Data Package $\mathcal Q$ to a Gauge-Fixed Path Integral for the
Trajectory-Induced $SU(3)$ Lattice Theory}, 2026.

\bibitem{os}
K.~Osterwalder and R.~Schrader,
\emph{Axioms for Euclidean Green's functions}, Commun. Math. Phys. 31 (1973),
83--112; and \emph{Axioms for Euclidean Green's functions. II}, Commun. Math.
Phys. 42 (1975), 281--305.

\end{thebibliography}

\end{document}
